\section{CLOUD\_MODE: arquitectura en la nube}
\label{sec:cloud_mode}

CLOUD\_MODE está compuesto por cuatro microservicios desplegados sobre
Railway, más un firmware ESP32-S3 que únicamente consume resultados.

\subsection{Servicio de autenticación}
El \textit{auth-service} es dueño de las cuentas de usuario y emite
tokens JWT firmados con RS256. Expone \texttt{POST /register},
\texttt{POST /login}, \texttt{POST /refresh} y \texttt{GET
/public-key}. Los tokens de acceso expiran a los 15 minutos y los de
refresco a los 30 días. El resto de los servicios valida los tokens
localmente contra la llave pública, sin consultar nunca la base de
usuarios.

\subsection{API Gateway}
El \textit{api-gateway} es el único servicio con el que la aplicación
móvil se comunica directamente. Valida el JWT, recibe la foto
(\texttt{POST /images}, \textit{multipart}), la sube a almacenamiento
tipo S3, invoca internamente al servicio de clasificación por HTTP,
guarda el resultado y lo devuelve a la aplicación. También expone
\texttt{GET /images} (historial) y \texttt{POST /devices/pair} para
emparejar un ESP32 con un usuario, generando sus credenciales MQTT.

\subsection{Servicio de clasificación}
El \textit{classification-service} ejecuta el algoritmo de
clasificación sobre la imagen recibida y devuelve
\texttt{\{label, confidence, stub, reason, features\}}. Publica el
mismo resultado en un tópico MQTT/TLS para que el ESP32 emparejado lo
reciba. Este servicio nunca es invocado directamente por la aplicación
móvil, sino siempre a través del API Gateway.

\subsection{Broker MQTT}
Eclipse Mosquitto, exclusivamente sobre TLS (puerto 8883), con listas
de control de acceso por dispositivo. Los resultados se publican en
\texttt{results/<usuario\_mqtt>}; un tópico
\texttt{commands/<...>} queda reservado para una futura actuación
física, aún no implementada.

\subsection{Firmware ESP32-S3 (visualización)}
Este dispositivo no toma fotos ni clasifica: se conecta por WiFi y por
MQTT sobre TLS, se suscribe a su propio tópico de resultados y muestra
cada clasificación entrante (opcionalmente en una pantalla OLED, o por
puerto serial).
